\section{智慧能源中的深度学习应用基础}

\subsection{典型任务}
从应用角度看，智慧能源中的深度学习任务可以概括为三类。第一类是预测类任务，如电力负荷预测、风电功率预测、光伏功率预测和建筑能耗预测；第二类是识别与诊断类任务，如故障识别、异常检测和设备健康评估；第三类是决策与控制类任务，如微电网能量管理、储能调度和需求响应优化。三类任务虽然形式不同，但都高度依赖数据质量、时间关联和场景适应性，因此非常适合使用数据驱动模型建模。

在本文聚焦的新能源预测场景中，深度学习的优势主要体现在三个方面。第一，新能源出力受辐照度、温度、云层变化、风速和风向等因素共同影响，呈现明显非线性，传统线性模型很难稳定描述这种关系。第二，功率序列本身存在较强时间依赖，LSTM 和 GRU 等循环结构更适合捕捉这种动态变化。第三，随着采集条件改善，模型输入已经不再局限于历史功率一个变量，而是可以同时包含多种气象与设备特征，多源数据融合成为进一步提升精度的重要方向。

\subsection{常见模型与评价指标}
目前智慧能源研究中最常见的深度学习模型包括 LSTM、GRU、CNN、TCN、Transformer 及其混合结构。LSTM 和 GRU 对长期依赖关系建模能力较强，因此在负荷与新能源预测中应用最广。CNN 可以提取局部模式，常与 LSTM 结合形成 CNN-LSTM，用于同时学习局部变化与时间依赖。TCN 通过膨胀卷积增强了长序列建模能力，在高频多步预测中表现突出。近年来，Transformer 和注意力机制也逐步进入新能源预测领域，用于处理更长时间跨度和更复杂的多变量关系。

预测模型通常采用 MAE、RMSE、MAPE 等指标进行评价。这三个指标的计算方式如下：

\[
\mathrm{MAE}=\frac{1}{n}\sum_{i=1}^{n}\left|y_i-\hat{y}_i\right|
\]

\[
\mathrm{RMSE}=\sqrt{\frac{1}{n}\sum_{i=1}^{n}\left(y_i-\hat{y}_i\right)^2}
\]

\[
\mathrm{MAPE}=\frac{100\%}{n}\sum_{i=1}^{n}\left|\frac{y_i-\hat{y}_i}{y_i}\right|
\]

其中，$y_i$ 表示真实值，$\hat{y}_i$ 表示预测值，$n$ 表示样本数。

MAE 便于理解平均偏差大小，RMSE 对大误差更敏感，更适合衡量极端天气或波动时段下的稳定性，MAPE 则便于跨场景比较，但在真实值接近零时容易失真，因此在夜间光伏出力场景下需要谨慎解释。对智慧能源而言，模型价值不能只看平均误差，还应结合鲁棒性和下游调度收益综合判断。

\newpage

\begin{center}
\captionof{table}{智慧能源典型任务、输入数据与常用模型\\Table 1 Typical tasks, inputs, and models in smart energy}
\vspace{3pt}
{\fontsize{8.6pt}{12pt}\selectfont
\setlength{\tabcolsep}{2pt}
\begin{tabular}{@{}p{0.20\columnwidth}p{0.28\columnwidth}p{0.23\columnwidth}p{0.21\columnwidth}@{}}
\toprule
任务 & 常见输入数据 & 常用模型 & 直接应用价值 \\
\midrule
电力负荷预测 & 历史负荷、气温、节假日、时段标签 & LSTM、GRU、Transformer、XGBoost & 支撑购售电计划与需求响应 \\
光伏功率预测 & 历史功率、辐照度、温度、湿度、云量 & LSTM、BiLSTM、CNN-LSTM、TCN & 提高并网消纳与储能控制效果 \\
风电功率预测 & SCADA、风速、风向、气压、数值天气预报 & LSTM、TCN、GNN、Attention & 降低风电波动带来的调度风险 \\
微电网能量管理 & 负荷、发电、储能状态、电价、运行约束 & DQN、DDPG、MPC+深度学习 & 降低运行成本并提升自治能力 \\
\bottomrule
\end{tabular}}
\label{tab:tasks-models}
\end{center}
